\section{Weighted comparison along a closed geodesic}
\label{weighted:section}

We retain $D=\pi/2$, $T=\pi$, $d=n-1$, and $m=d-k$, with
$m,k\geq1$. The factorization $f=f_Hf_K$ turns the logarithm of
the ratio to the model density into a sum of two block contributions.
We will show that each contribution has nonnegative integral with
weight $w(|y-x|)$. The scalar identity below cancels the integrated
endpoint terms in the block bounds. A strict-convexity remainder
will then give the determinant comparison and identify its equality
case, without separately classifying equality in the matrix estimates.

The scalar identity is a logarithmic form of the weighted cancellation
in Kazdan's proof of the Berger--Kazdan inequality
\citep[Appendix~E, especially E.4--E.8]{Kazdan1978AppendixE}.
At interval length $\pi$, writing Kazdan's profile as
$\phi(t)=\sin(t)e^{u(t)}$ gives $p=\phi^{-2}$,
$p_0(t)=\sin^{-2}t$, $g=\log(p/p_0)=-2u$, and weight
$\rho(t)=W(t)/\sin t$.
We give the version needed here, with bounded logarithmic ratios
and a measurable weight, including the endpoint integrability.
Its application to the two factors depends on the folded symmetry
proved below.

\begin{lemma}[Weighted scalar identity]\label{weighted:scalar}
Let $b>0$, $\lambda=\pi/b$, and
\[
 p_0(t)=\frac{\lambda^2}{\sin^2(\lambda t)},\qquad 0<t<b.
\]
Let $p:(0,b)\to(0,\infty)$ be continuous and suppose that
$g=\log(p/p_0)$ is bounded.  Let $W$ be nonnegative and measurable on
$(0,b)$, with $W(b-z)=W(z)$ almost everywhere and
$W(z)/\sin(\lambda z)$ essentially bounded.  Write
\[
 \Delta_b=\{(x,y):0<x<y<b\},\qquad
 d\nu_{xy}(t)=\frac{p_0(t)\,dt}{\int_x^y p_0},\quad x<t<y.
\]
Then all the following integrals are finite, and
\begin{align}
 &\int_{\Delta_b}W(y-x)
 \left[\log\frac{\int_x^y p}{\int_x^y p_0}
       -\frac{g(x)+g(y)}2\right]dx\,dy \notag\\
 &\qquad=\int_{\Delta_b}W(y-x)
 \left[\log\int_x^y e^{g(t)}\,d\nu_{xy}(t)
       -\int_x^y g(t)\,d\nu_{xy}(t)\right]dx\,dy
 \ \geq\ 0.\label{weighted:scalar-identity}
\end{align}
\end{lemma}

\begin{proof}
Direct integration gives
\begin{equation}
 P_0(x,y):=\int_x^y p_0
 =\frac{\lambda\sin(\lambda(y-x))}
        {\sin(\lambda x)\sin(\lambda y)},\qquad
 \frac{P_0(x,y)}{\sqrt{p_0(x)p_0(y)}}
 =\frac{\sin(\lambda(y-x))}{\lambda}.
 \label{weighted:chord}
\end{equation}
The measure $\nu_{xy}$ is a probability measure.  Both its average of
$g$ and $\log\int e^g\,d\nu_{xy}=\log(\int p/\int p_0)$ have absolute
value at most $\|g\|_\infty$.  Since $W$ is essentially bounded, the asserted
finiteness follows already from these bounds.  Jensen's inequality makes
the integrand on the second line of
\eqref{weighted:scalar-identity} nonnegative.  It remains to establish
the equality of the integrated linear terms.

Set $K(z)=W(z)/\sin(\lambda z)$ and $I_W=\int_0^b W(z)\,dz$, and define
\[
 F(t)=\int_{0<x<t<y<b}
 K(y-x)\sin(\lambda x)\sin(\lambda y)\,dx\,dy.
\]
Here the two-variable kernel
$q(x,y)=K(y-x)\sin(\lambda x)\sin(\lambda y)$ is essentially bounded.
Fubini gives
\[
 F(t)=\int_0^t\left(\int_x^b q(x,y)\,dy\right)dx
       -\int_0^t\left(\int_0^y q(x,y)\,dx\right)dy.
\]
Thus $F$ is absolutely continuous on $[0,b]$, $F(0)=0$, and, for almost
every $t$,
\begin{align*}
 F'(t)
 &=\sin(\lambda t)\left[
   \int_0^{b-t}K(z)\sin(\lambda(t+z))\,dz
   -\int_0^t K(z)\sin(\lambda(t-z))\,dz\right]\\
 &=\sin(\lambda t)\int_0^b K(z)\sin(\lambda(t+z))\,dz\\
 &=I_W\sin(\lambda t)\cos(\lambda t).
\end{align*}
For the second equality substitute $b-z$ in the subtracted integral:
$K(b-z)=K(z)$ and
$\sin(\lambda(t-b+z))=-\sin(\lambda(t+z))$.
For the last equality, symmetry gives
$\int_0^b K(z)\cos(\lambda z)\,dz=0$, while
$\int_0^b K(z)\sin(\lambda z)\,dz=I_W$.
Integration therefore yields
\begin{equation}
 F(t)=\frac{I_W}{2\lambda}\sin^2(\lambda t),\qquad
 \frac{p_0(t)}{\lambda}F(t)=\frac{I_W}{2}.
 \label{weighted:coefficient}
\end{equation}
This argument has not differentiated the measurable weight.

Using \eqref{weighted:chord} and then
\eqref{weighted:coefficient}, the integrated interval averages are
\begin{align*}
 \int_{\Delta_b}W(y-x)\int_x^y g(t)\,d\nu_{xy}(t)\,dx\,dy
 &=\int_0^b g(t)\frac{p_0(t)}{\lambda}F(t)\,dt\\
 &=\frac{I_W}{2}\int_0^b g(t)\,dt.
\end{align*}
The exchange is absolutely integrable: replacing $g$ by $|g|$ in the
left side bounds it by
$\|g\|_\infty\int_{\Delta_b}W(y-x)\,dx\,dy$.
In particular, no integral of $p_0$ over the entire interval is taken.
The endpoint terms have exactly the same value, since
\begin{align*}
 &\frac12\int_{\Delta_b}W(y-x)(g(x)+g(y))\,dx\,dy\\
 &\qquad=\frac12\int_0^b g(t)
 \left[\int_0^{b-t}W(z)\,dz+\int_0^t W(z)\,dz\right]dt
 =\frac{I_W}{2}\int_0^b g(t)\,dt.
\end{align*}
The final equality again uses $W(b-z)=W(z)$.  Subtracting these equal
finite quantities proves the identity and hence the lemma.
\end{proof}

\begin{theorem}[Comparison on each geodesic]\label{weighted:geodesic}
In the normalized intermediate case of Proposition~\ref{geom:structure},
let $f(x,y)$ be the Jacobi determinant of any fixed unit-speed geodesic,
as in Proposition~\ref{lag:determinant}.  Put
\[
 \mathcal R=(0,T)^2,\qquad
 w(z)=(\sin z)^m\left|\frac{\sin(2z)}2\right|^k
 \quad(0\leq z\leq T).
\]
Then
\begin{equation}
 \int_{\mathcal R}f(x,y)\,dx\,dy
 \ \geq\ \int_{\mathcal R}w(|y-x|)\,dx\,dy.
 \label{weighted:comparison}
\end{equation}
Equality holds if and only if
$f(x,y)=w(|y-x|)$ for every $(x,y)\in[0,T]^2$.
\end{theorem}

\begin{proof}
All constants in this proof may depend on the chosen geodesic.
Write $w_{xy}=w(|y-x|)$.  We first work off the null set
\[
 E=\{(x,y)\in\mathcal R:
       x=D\text{ or }y=D\text{ or }x=y\text{ or }|y-x|=D\}.
\]
There $f,w_{xy}$, and the two determinant factors $f_H,f_K$ of
Proposition~\ref{lag:block-bounds} are positive, and $f=f_Hf_K$.  Define
\[
 L_H(x,y)=\log\frac{f_H(x,y)}{\sin^m(|y-x|)},\qquad
 L_K(x,y)=\log\frac{f_K(x,y)}{|\sin(2|y-x|)/2|^k}.
\]
Set their values on $E$ equal to zero for integration.
Both are symmetric in the actual endpoint slots $x,y$.

\paragraph{Finiteness of the logarithmic integrals.}
The profile bounds in Proposition~\ref{lag:profiles} give constants
$0<\alpha\leq\beta<\infty$ such that
\begin{align*}
 \alpha&\leq\frac{h(t)}{\csc^2t}\leq\beta
 &&(0<t<T),\\
 \alpha&\leq\frac{s_i(t)}{4\csc^2(2t)}\leq\beta
 &&(0<t<D,\ i=0,1).
\end{align*}
Set $\eta=\alpha/\beta\leq1$.
For $x<y$, the first bound in Proposition~\ref{lag:block-bounds} and the chord
identity \eqref{weighted:chord} imply
\[
 f_H(x,y)^{1/m}
 \geq\frac{\int_x^y h}{\sqrt{h(x)h(y)}}
 \geq\eta\sin(y-x).
\]
Symmetry gives this lower bound for either endpoint order.
Similarly, put $B_0(z)=\sin(2z)/2$ for $0<z<D$.
For $0<u<v<D$ and $i,j\in\{0,1\}$, the second block bound gives
\[
 f_K(u+iD,v+jD)^{1/k}
 \geq\frac{(\int_u^v s_i)^{1/2}(\int_u^v s_j)^{1/2}}
              {\sqrt{s_i(u)s_j(v)}}
 \geq\eta B_0(v-u).
\]
In these coordinates
$B_0(v-u)=|\sin(2|y-x|)/2|$.  Every pair outside $E$ has this form,
possibly after exchanging endpoints.  Consequently
\begin{equation}
 L_H\geq m\log\eta,
 \qquad L_K\geq k\log\eta,
 \qquad L_H+L_K=L:=\log(f/w_{xy}).
 \label{weighted:log-lower-bounds}
\end{equation}
With $a^+=\max(a,0)$, the elementary estimate
\[
 w_{xy}L^+\leq f
\]
and continuity of $f$ on the compact square prove integrability of
$w_{xy}L^+$.  The individual negative parts are bounded by
\eqref{weighted:log-lower-bounds}; also
\[
 L_H^+\leq L^+-k\log\eta,\qquad
 L_K^+\leq L^+-m\log\eta.
\]
Thus $w_{xy}L_H$, $w_{xy}L_K$, and $w_{xy}L$ are all absolutely
integrable.  The following additions of logarithmic integrals therefore
involve only finite quantities.

\paragraph{The $H$ factor on the full square.}
Apply Lemma~\ref{weighted:scalar} with
$b=T$, $\lambda=1$, $p=h$, and $W=w$.  Its assumptions hold by
Proposition~\ref{lag:profiles} and
\[
 w(T-z)=w(z),\qquad
 \frac{w(z)}{\sin z}
 =\sin^{m-1}z\left|\frac{\sin(2z)}2\right|^k,
\]
which is bounded because $m\geq1$.
Writing $g_H=\log(h/\csc^2t)$, the first block bound and
\eqref{weighted:chord} give, for $0<x<y<T$ outside $E$,
\[
 \frac{L_H(x,y)}m
 \geq\log\frac{\int_x^y h(t)\,dt}{\int_x^y\csc^2t\,dt}
       -\frac{g_H(x)+g_H(y)}2.
\]
After multiplication by $w(y-x)$, Lemma~\ref{weighted:scalar}
makes the integral of the right side nonnegative.  Multiplication by
$m$ and endpoint symmetry then yield
\begin{equation}
 \int_{\mathcal R}w_{xy}L_H(x,y)\,dx\,dy\geq0.
 \label{weighted:H-log}
\end{equation}

\paragraph{The $K$ factor and the four folded weights.}
Put $p_0(t)=4\csc^2(2t)$ on $(0,D)$ and define
\[
 g_i(t)=\log\frac{s_i(t)}{p_0(t)},\qquad
 A_i(u,v)=\log\frac{\int_u^v s_i(t)\,dt}{\int_u^v p_0(t)\,dt}
 \quad(i=0,1).
\]
For $(u,v)\in\Delta_D$, $z=v-u$, $x=u+iD$, and $y=v+jD$, the second
block bound gives the precise logarithmic inequality
\begin{equation}
 \frac{L_K(u+iD,v+jD)}k
 \geq\frac12\bigl[A_i(u,v)+A_j(u,v)-g_i(u)-g_j(v)\bigr].
 \label{weighted:folded-bound}
\end{equation}
Indeed, \eqref{weighted:chord} with $\lambda=2$ converts the reference
interval integral and its endpoint factors into $B_0(z)$.
The actual separation is $z$ if $i=j$, $D+z$ if $(i,j)=(0,1)$,
and $D-z$ if $(i,j)=(1,0)$.  In the last case $x>y$; the bound in
Proposition~\ref{lag:block-bounds} includes this order.
All three separations have the same absolute fast model factor
$B_0(z)$.

The four weights $w_{ij}(z)=w(|z+(j-i)D|)$ form the matrix
\begin{equation}
 (w_{ij}(z))_{i,j=0}^1
 =B_0(z)^k
 \begin{pmatrix}
  \sin^m z&\cos^m z\\
  \cos^m z&\sin^m z
 \end{pmatrix}.
 \label{weighted:weight-matrix}
\end{equation}
Every row and every column sums to
\[
 W_D(z)=(\sin^m z+\cos^m z)B_0(z)^k.
\]
This is symmetric under $z\mapsto D-z$, and
\[
 \frac{W_D(z)}{\sin(2z)}
 =2^{-k}(\sin^m z+\cos^m z)\sin^{k-1}(2z)
\]
is bounded for every $k\geq1$.  When $k=1$, its limit at either endpoint
is $1/2$, so this case also satisfies Lemma~\ref{weighted:scalar}.

Multiply \eqref{weighted:folded-bound} by $w_{ij}(z)$ and sum over
$i,j$.  The row sums control the terms indexed by $i$, and the column
sums control those indexed by $j$, giving exactly
\begin{align}
 \sum_{i,j=0}^1 w_{ij}(z)\frac{L_K(u+iD,v+jD)}k
 &\geq W_D(z)\sum_{i=0}^1
 \left[A_i(u,v)-\frac{g_i(u)+g_i(v)}2\right].
 \label{weighted:folded-sum}
\end{align}
Lemma~\ref{weighted:scalar}, applied on $(0,D)$ separately to $s_0$
and $s_1$ with weight $W_D$, shows that the right side has nonnegative
integral over $\Delta_D$.

For completeness, the four images of $\Delta_D$ used above consist
precisely of pairs whose first endpoint has smaller residue modulo $D$
than the second.  Their transposes contain precisely the opposite
residue order.  These eight disjoint regions cover $\mathcal R\setminus E$;
each coordinate change has Jacobian one.  Since $w_{xy}L_K(x,y)$ is
symmetric, it follows that
\begin{align}
 &\int_{\mathcal R}w_{xy}L_K(x,y)\,dx\,dy\notag\\
 &\qquad=2k\int_{\Delta_D}\sum_{i,j=0}^1
 w_{ij}(v-u)\frac{L_K(u+iD,v+jD)}k\,du\,dv
 \ \geq\ 0.
 \label{weighted:K-log}
\end{align}
Here the factor $2$ accounts for endpoint exchange and the factor $k$
restores the logarithm of $f_K$ from that of its $k$th root.

\paragraph{Comparison and equality.}
Adding \eqref{weighted:H-log} and \eqref{weighted:K-log} gives
\[
 \int_{\mathcal R}w_{xy}\log(f/w_{xy})\,dx\,dy\geq0.
\]
For $r>0$, the strictly convex function
\[
 \Phi(r)=r-1-\log r
\]
is nonnegative and vanishes exactly at $r=1$.  On
$\mathcal R\setminus E$ put $r=f/w_{xy}$.  All terms in the identity
\begin{equation}
 \int_{\mathcal R}(f-w_{xy})\,dx\,dy
 =\int_{\mathcal R}w_{xy}\log r\,dx\,dy
  +\int_{\mathcal R}w_{xy}\Phi(r)\,dx\,dy
 \label{weighted:strict-gap}
\end{equation}
are finite by the preceding integrability argument.  Both terms on
the right are nonnegative, proving \eqref{weighted:comparison}.
If equality holds, the last integral is zero.  Since $w_{xy}>0$ off
the diagonal and cut lines, $\Phi(r)=0$ almost everywhere, hence
$f=w_{xy}$ almost everywhere in $\mathcal R$.
Continuity of both functions on $[0,T]^2$ extends this identity to
every pair, including the exceptional lines and boundary.  Conversely,
that pointwise identity gives equality in \eqref{weighted:comparison}.
\end{proof}
